\begin{minipage}{0.7\textwidth}
La lame d’une scie sauteuse effectue un MHS vertical avec une fréquence de $\v{Freq}$. La distance entre la position la plus haute et la plus basse de la lame est de $\v{Course}$. Lorsqu’elle se trouve à sa position la plus basse, la lame dépasse de $\v{Depassement}$ sous la planche.

\begin{enumerate}[label=\alph*)]
    \item Écrivez l’équation de la position de la lame en fonction du temps (vous pouvez supposer que la constante de phase est nulle).
    \item Au cours d’un cycle de l’oscillation, pendant combien de temps la lame dépasse-t-elle de la planche ? Tracez un graphique approximatif de la position en fonction du temps et identifiez sur votre graphique les points utilisés pour le calcul. Donnez votre réponse en millisecondes.
\end{enumerate}

\end{minipage}
\hfill
\begin{minipage}{0.25\textwidth}
\centering
\includegraphics[width=1 \linewidth]{Figures/scie-sauteuse.png}
\end{minipage}
